% fbd_incline: free-body diagram of a block at rest on a rough incline (Ch 7, fig:fbd)
% theta = 30 degrees; forces to scale for equilibrium:
%   mg = 2.4 cm, N = mg cos(theta) = 2.08 cm, f = mg sin(theta) = 1.2 cm (up the slope,
%   since the block tends to slide down).
% N and f act at the contact (centre of the base), mg at the centre of mass.
% mg is resolved into a right triangle: mg sin(theta) along the slope (down) and
% mg cos(theta) into the slope; the angle between mg and the second leg is theta.
\begin{tikzpicture}[x=1cm,y=1cm]
  \def\th{30}
  \def\Lb{8.2}
  \def\W{2.4}
  \def\s{5.5}
  \pgfmathsetmacro\Hh{\Lb*tan(\th)}
  \pgfmathsetmacro\Nl{\W*cos(\th)}
  \pgfmathsetmacro\fl{\W*sin(\th)}
  % incline and floor
  \path[ground,draw=none] (-0.35,-0.14) rectangle (\Lb+0.35,0);
  \fill[cSlope] (0,0) -- (\Lb,0) -- (\Lb,\Hh) -- cycle;
  \draw[outline] (-0.35,0) -- (\Lb+0.35,0);
  \draw[outline] (0,0) -- (\Lb,\Hh) -- (\Lb,0);
  \draw[angle arc] (1.15,0) arc[start angle=0,end angle=\th,radius=1.15];
  \node at (\th/2:1.5) {$\theta$};
  % block
  \coordinate (P) at (\th:\s);
  \begin{scope}[shift={(P)},rotate=\th]
    \draw[body] (-0.75,0) rectangle (0.75,0.95);
    \coordinate (C) at (0,0.475);
  \end{scope}
  % components of mg (drawn first, so the forces sit on top)
  \coordinate (T) at ($(C)+(0,-\W)$);
  \coordinate (A) at ($(C)+(180+\th:\fl)$);
  \draw[component,cCrim] (C) -- (A)
    node[anchor=south east,inner sep=1pt] {$mg\sin\theta$};
  \draw[component,cCrim] (A) -- (T)
    node[pos=0.86,anchor=east,inner sep=3pt] {$mg\cos\theta$};
  \draw[angle arc,cCrim] ($(T)+(90:0.5)$) arc[start angle=90,end angle=90+\th,radius=0.5];
  \node[cCrim] at ($(T)+(90+\th/2:0.78)$) {$\theta$};
  % forces
  \draw[force] (C) -- (T) node[anchor=north] {$mg$};
  \draw[force] (P) -- ++(90+\th:\Nl) node[anchor=south east,inner sep=1pt] {$N$};
  \draw[applied] (P) -- ++(\th:\fl) node[anchor=north west,inner sep=2pt] {$f$};
  \node[contact point] at (P) {};
\end{tikzpicture}
